\newpage

\section{Simulation Studies and Real Data Analysis}

\subsection{Simulation Studies}

We conducted extensive simulation studies to evaluate the finite-sample performance of the proposed penalized spline estimator for generalized partially linear single-index models (GPLSIMs). The objectives were to assess (i) the estimation accuracy of the regression and index parameters $(\Bb, \Ba)$, (ii) the recovery of the nonlinear link function $\vp(\cdot)$, and (iii) the performance of several variance estimators. 

% Simulation Design and Estimation Procedure
Two canonical generalized linear model settings were considered: a binomial model with a logistic link and a Poisson model with a log link. In both cases, data were generated from the $g\{\mu_i\} = \Bb^\top \bx_i + \vp(\Ba^\top \bz_i)$, where $\Ba \in \mathbb{R}^p$ is the index parameter, $\Bb \in \mathbb{R}^d$ is the vector of regression coefficients, and $\vp(\cdot)$ is an unknown smooth function. The true parameters were specified as
$\Ba_0 = \frac{1}{\sqrt{3}}(1,1,1)^\top, 
\Bb_0 = (1, -0.5)^\top,$
and the nonlinear function was defined as $\vp(u) = \sin(2u)$. To ensure identifiability, the nonlinear component was empirically centered by subtracting $\mathbb{E}\{\vp(\Ba^\top \bz_i)\}$.

The covariates were generated independently as $Z_1, Z_2, Z_3 \sim \text{Uniform}(0,2)$,   $X_1 \sim N(0,1)$, $X_2 \sim \text{Bernoulli}(0.5)$. We denote $  \bz_i = (Z_{i1}, Z_{i2},Z_{i3})^\top$ and $\bx_i = (X_{i1}, X_{i2})^\top$. Poisson outcomes are generated by $Y_i \sim \text{Poisson}(\mu_i)$, $\mu_i =  C_i\exp\left(\Bb^\top \bx_i + \vp(\Ba^\top \bz_i) \right)$, with an offset $C_i \sim \text{Uniform}(0,10)$. Binomial outcomes are generated by $Y_i \sim \text{Bernoulli}(\mu_i)$, $\mu_i = \frac{\exp(\eta_i)}{1 + \exp(\eta_i)}$, with $\eta_i = \Bb^\top \bx_i + \vp(\Ba^\top \bz_i)$. The sample size was either $n=400$, $n=800$ or $n=2000$ with $1000$ replications. The unknown function $\vp(\cdot)$ was approximated using a cubic B-spline basis with a second-order difference penalty. Parameters were estimated jointly via an iterative two-stage algorithm consisting of Fisher scoring with step-halving and smoothing parameter updates using the generalized Fellner–Schall method. The index parameter $\Ba$ was recovered from the reparameterized vector $\Be$ through $\Ba = \frac{(1, \Be^\top)^\top}{\sqrt{1 + \|\Be\|^2}}$

\begin{table}[ht!]
\centering
\begin{threeparttable}
\caption{Summary of parameter estimation for Poisson and Binomial simulation studies.}
\label{tab:sim_param_pois_bino_main}
\setlength{\tabcolsep}{3.5pt}
\begin{tabular}{llcccccccc}
\toprule
& & \multicolumn{4}{c}{Poisson} & \multicolumn{4}{c}{Binomial} \\
\cmidrule(lr){3-6} \cmidrule(lr){7-10}
& & \multicolumn{2}{c}{$n=400$} & \multicolumn{2}{c}{$n=800$}
  & \multicolumn{2}{c}{$n=400$} & \multicolumn{2}{c}{$n=800$} \\
\cmidrule(lr){3-4} \cmidrule(lr){5-6} \cmidrule(lr){7-8} \cmidrule(lr){9-10}
Parameter & Metric
& Lu & Yu & Lu & Yu
& Lu & Yu & Lu & Yu \\

\midrule

\multirow{4}{*}{$\alpha_{1}$}
& Mean & 0.5767 & 0.5624 & 0.5774 & 0.5734 & 0.5610 & 0.5376 & 0.5724 & 0.5595 \\
& Bias & 0.0006 & 0.0149 & $-0.0001$ & 0.0040 & 0.0164 & 0.0397 & 0.0049 & 0.0178 \\
& SD   & 0.0176 & 0.1154 & 0.0123 & 0.0431 & 0.1335 & 0.2183 & 0.0834 & 0.1522 \\
& MSE  & 0.0003 & 0.0135 & 0.0002 & 0.0019 & 0.0181 & 0.0492 & 0.0070 & 0.0234 \\
\midrule

\multirow{4}{*}{$\alpha_{2}$}
& Mean & 0.5774 & 0.5430 & 0.5777 & 0.5787 & 0.5610 & 0.3999 & 0.5704 & 0.5057 \\
& Bias & $-0.0000$ & 0.0343 & $-0.0003$ & $-0.0013$ & 0.0163 & 0.1775 & 0.0070 & 0.0716 \\
& SD   & 0.0178 & 0.1925 & 0.0125 & 0.0437 & 0.1389 & 0.4216 & 0.0839 & 0.2878 \\
& MSE  & 0.0003 & 0.0382 & 0.0002 & 0.0019 & 0.0195 & 0.2091 & 0.0071 & 0.0879 \\
\midrule

\multirow{4}{*}{$\alpha_{3}$}
& Mean & 0.5771 & 0.5433 & 0.5765 & 0.5751 & 0.5625 & 0.4059 & 0.5709 & 0.5094 \\
& Bias & 0.0003 & 0.0340 & 0.0008 & 0.0023 & 0.0149 & 0.1715 & 0.0065 & 0.0679 \\
& SD   & 0.0180 & 0.2081 & 0.0123 & 0.0438 & 0.1310 & 0.4016 & 0.0840 & 0.2566 \\
& MSE  & 0.0003 & 0.0444 & 0.0002 & 0.0019 & 0.0174 & 0.1905 & 0.0071 & 0.0704 \\
\midrule

\multirow{4}{*}{$\beta_{1}$}
& Mean & 1.0003 & 1.0026 & 0.9996 & 0.9972 & 1.0170 & 1.0186 & 1.0094 & 1.0072 \\
& Bias & $-0.0003$ & $-0.0026$ & 0.0004 & 0.0028 & $-0.0170$ & $-0.0186$ & $-0.0094$ & $-0.0072$ \\
& SD   & 0.0166 & 0.0716 & 0.0111 & 0.0520 & 0.1422 & 0.1494 & 0.0985 & 0.0999 \\
& MSE  & 0.0003 & 0.0051 & 0.0001 & 0.0027 & 0.0205 & 0.0226 & 0.0098 & 0.0100 \\
\midrule

\multirow{4}{*}{$\beta_{2}$}
& Mean & $-0.5005$ & $-0.5055$ & $-0.4996$ & $-0.5022$ & $-0.5031$ & $-0.5096$ & $-0.5003$ & $-0.5004$ \\
& Bias & 0.0005 & 0.0055 & $-0.0004$ & 0.0022 & 0.0031 & 0.0096 & 0.0003 & 0.0004 \\
& SD   & 0.0346 & 0.1171 & 0.0235 & 0.0821 & 0.1647 & 0.2364 & 0.1146 & 0.1594 \\
& MSE  & 0.0012 & 0.0137 & 0.0006 & 0.0067 & 0.0271 & 0.0559 & 0.0131 & 0.0254 \\

\bottomrule
\end{tabular}

\begin{tablenotes}
\footnotesize
\item $\boldsymbol{\alpha}_0=(1/\sqrt{3},1/\sqrt{3},1/\sqrt{3})^\top$ and 
$\boldsymbol{\beta}_0=(1,-0.5)^\top$. Mean, bias, standard deviation (SD), and mean squared error (MSE) are based on 1000 Monte Carlo replications.
\end{tablenotes}
\end{threeparttable}
\end{table}

% Parameter Estimation Results
Tables~\ref{tab:sim_param_pois_bino_main} summarizes the parameter estimation results. In the Poisson model, the proposed estimator (Lu) exhibits negligible bias and substantially smaller variability than the competing method~\cite{Yu2017} (Yu). For example, at $n=800$, the standard deviation of $\alpha_1$ is reduced from 0.0431 under Yu to 0.0123 under Lu, corresponding to a reduction of approximately 71\%. Similar efficiency gains are observed for the other index parameters, with Lu producing much smaller SD and MSE values across $\alpha_1$, $\alpha_2$, and $\alpha_3$. The parametric components are also accurately recovered, with very small bias for both $\beta_1$ and $\beta_2$.
In the Binomial model, estimation is more challenging, especially for the index parameters. Nevertheless, Lu remains more stable than Yu. At $n=800$, Yu shows noticeable bias for $\alpha_2$ and $\alpha_3$ with biases of 0.0716 and 0.0679, respectively, whereas Lu reduces these to 0.0070 and 0.0065. Lu also has considerably smaller variability for the index parameters; for example, the SD of $\alpha_2$ decreases from 0.2878 under Yu to 0.0839 under Lu. As the sample size increases to $n=2000$ (See Tables~\ref{tab:sim_param_pois_extended}--\ref{tab:sim_param_bino_extended}), both methods improve, but Lu continues to show smaller SD and MSE values for most parameters. Overall, the results indicate that the proposed estimator provides accurate and stable recovery of both the index and parametric components, with the strongest efficiency advantages observed for the index parameters.

% Tables~\ref{tab:sim_param_pois_bino_main}  summarize estimation results. In the Poisson model, the proposed estimator (Lu) exhibits negligible bias and substantially smaller variability than the competing method\cite{Yu2017} (Yu). For example, the standard deviation of $\alpha_1$ is reduced from approximately 0.043 to 0.012 at $n=800$, representing a reduction of over 70\%. Similar gains are observed across all index parameters. In the Binomial model, estimation is more challenging due to the nonlinear link. Nevertheless, the proposed estimator maintains stable performance with smaller bias and variance, particularly for the index parameters. At $n=800$, the competing method shows substantial bias (e.g., $\alpha_2$: 0.0716 vs 0.0070), whereas the proposed estimator remains close to the truth. Performance improves for both methods at $n=2000$, but the proposed estimator consistently retains an efficiency advantage. These results demonstrate accurate and stable recovery of both parametric and index components.

% Tables~\ref{tab:sim_param_pois} and \ref{tab:sim_param_bino} summarize the estimation performance for the Poisson and Binomial models, respectively. Across both models, the proposed estimator (Lu) demonstrates substantially improved efficiency compared to the competing method \cite{Yu2017} (Yu). In the Poisson setting, the proposed estimator achieves extremely small bias and variance for all parameters, with standard deviations decreasing markedly as the sample size increases from $n=800$ to $n=2000$. The gains in efficiency are particularly notable for the index parameters $(\alpha_1,\alpha_2,\alpha_3)$, where the variability of the proposed estimator is consistently much smaller than that of the competing method. In the Binomial setting, estimation is inherently more challenging due to the nonlinear link and bounded outcome. Nevertheless, the proposed estimator maintains stable and accurate performance, with moderate bias that diminishes as the sample size increases. In contrast, the competing method exhibits noticeably larger bias and variability, particularly for the index parameters at $n=800$. The improvement at $n=2000$ confirms the consistency and asymptotic efficiency of the proposed estimator. Overall, these results indicate that the proposed method provides reliable and precise estimation of both parametric and index components, even under nonlinear and non-Gaussian settings.

% \begin{figure}[ht!]
%     \centering

%     %---- First Row (N=400)----%
%     \begin{subfigure}[t]{0.48\textwidth}
%         \centering
%         \includegraphics[width=\textwidth]{Poisson_phi_estimation_plot400_260423.pdf}
%         \caption{Poisson model ($n=400$)}
%         \label{fig:phi_pois_800}
%     \end{subfigure}
%     \hfill
%     \begin{subfigure}[t]{0.48\textwidth}
%         \centering
%         \includegraphics[width=\textwidth]{Binomial_phi_estimation_plot400_260423.pdf}
%         \caption{Binomial model ($n=400$)}
%         \label{fig:phi_binom_800}
%     \end{subfigure}

%     %---- Second Row (N=800)----%
%     \begin{subfigure}[t]{0.48\textwidth}
%         \centering
%         \includegraphics[width=\textwidth]{Poisson_phi_estimation_plot800_260423.pdf}
%         \caption{Poisson model ($n=800$)}
%         \label{fig:phi_pois_800}
%     \end{subfigure}
%     \hfill
%     \begin{subfigure}[t]{0.48\textwidth}
%         \centering
%         \includegraphics[width=\textwidth]{Binomial_phi_estimation_plot800_260423.pdf}
%         \caption{Binomial model ($n=800$)}
%         \label{fig:phi_binom_800}
%     \end{subfigure}
    
%     % %---- Third Row (N=2000) ----%
%     % \begin{subfigure}[t]{0.48\textwidth}
%     %     \centering
%     %     \includegraphics[width=\textwidth]{Poisson_phi_estimation_plot2000_260423.pdf}
%     %     \caption{Poisson model ($n=2000$)}
%     %     \label{fig:phi_pois_2000}
%     % \end{subfigure}
%     % \hfill 
%     % \begin{subfigure}[t]{0.48\textwidth}
%     %     \centering
%     %     \includegraphics[width=\textwidth]{Binomial_phi_estimation_plot2000_260423.pdf}
%     %     \caption{Binomial model ($n=2000$)}
%     %     \label{fig:phi_binom_2000}
%     % \end{subfigure}
%     \vspace{0.5cm}
       
%     \caption{\textbf{Estimated Link Function.} 
%     Comparison of the true link function with the mean estimated link function and its pointwise 95\% confidence bands for Poisson and Binomial models, based on 1,000 Monte Carlo replications. The left panels correspond to the Poisson model, while the right panels correspond to the Binomial model. %Increasing the sample size from $n=400$ to $n=800$ results in improved estimation accuracy and narrower confidence bands, illustrating the consistency and efficiency of the proposed method.
%     }
    
%     \label{fig:estimated_link_function}
% \end{figure}

% \newpage

%Variance Estimation and Evaluation Metrics
% Next, four standard deviation estimators for $(\Bb, \Be)$ were compared: Model-based estimator (\texttt{se.mod}); Robust sandwich estimator (\texttt{se.rob}); Leverage-adjusted robust estimator (\texttt{se.lev}); and Finite-sample corrected estimator (\texttt{se.fin}). Performance was evaluated using bias, standard deviation (SD), root mean squared error (RMSE), and empirical coverage probabilities (CP) of 95\% Wald confidence intervals. The accuracy of the estimated link function was assessed through visual comparison and {\color{red} integrated squared error (ISE)}.

\begin{figure}[ht!]
\centering
\includegraphics[width=0.95\textwidth]{variance_summary_two_panel.png}
\caption{
Summary of variance estimation performance for the Poisson and Binomial simulation studies, based on 1,000 Monte Carlo replications.
Panel (A) displays the ratio of the model-based standard error to the empirical standard deviation, where the dashed horizontal line at 1 indicates perfect calibration.
Panel (B) displays the empirical coverage probability of the nominal 95\% confidence intervals, where the dashed horizontal line indicates the target coverage level of 0.95.
Red curves correspond to the $\alpha$ parameters and blue curves correspond to the $\beta$ parameters.
}
\label{fig:variance_summary_main}
\end{figure}

%Variance Estimation Results
We compared four variance estimators: model-based (\texttt{se.mod}), robust (\texttt{se.rob}), leverage-adjusted (\texttt{se.lev}), and finite-sample corrected (\texttt{se.fin}). Results are reported in Tables~\ref{tab:variance_study_pois} and \ref{tab:variance_study_bino}. In the Poisson model, all estimators perform similarly, with empirical standard deviations closely matching estimated standard errors and coverage probabilities near the nominal 95\% level. In contrast, the Binomial model shows systematic underestimation of variability by the model-based estimator at small ($n=400$) or moderate ($n=800$) sample sizes (See Table~\ref{tab:variance_study_bino}). 
%at $n=800$, leading to undercoverage (approximately 0.91–0.92 for index parameters). 
The corrected estimators improve coverage, with the leverage-adjusted correction providing the most stable performance. At $n=2000$, all estimators converge, and coverage approaches the nominal level. These results indicate that variance correction is necessary for valid inference in moderate sample sizes, particularly for nonlinear models.

%Variance Estimation Results
% Tables~\ref{tab:variance_study_pois} and \ref{tab:variance_study_bino} report the performance of four variance estimators.
% In the Poisson model, all variance estimators perform similarly, with empirical standard deviations closely matching the estimated standard errors. Coverage probabilities are near the nominal 95\% level across all parameters and sample sizes, indicating that the model-based estimator is adequate in this setting. In contrast, the Binomial model reveals more nuanced behavior. The model-based estimator tends to underestimate variability, particularly for the index parameters at $n=800$, leading to undercoverage (approximately 0.91–0.92). This reflects the increased curvature and nonlinearity of the likelihood in the Binomial case. The robust and leverage-adjusted estimators provide slight improvements, while the finite-sample corrected estimator achieves the most stable performance overall. As the sample size increases to $n=2000$, all estimators converge and coverage probabilities approach the nominal level. These findings highlight that variance correction can be effective in moderate sample sizes, especially for nonlinear models, while asymptotic approximations become adequate in larger samples.
%Variance Estimation Results
Figure~\ref{fig:variance_summary_main} summarizes the finite-sample performance of the proposed variance estimation procedures for the Poisson and Binomial GPLSIMs across different sample sizes. Panel~(A) compares the ratio of the model-based standard error to the empirical standard deviation, where values close to one indicate accurate variance estimation. For the Poisson model, the ratios remain close to one for both the index parameters $\alpha$ and the regression parameters $\beta$ across all sample sizes, indicating excellent calibration of the model-based variance estimator. In contrast, the Binomial model exhibits noticeable underestimation of variability for the index parameters when the sample size is small, with the ratios substantially below one at $n=400$ and $n=800$, although calibration improves as the sample size increases to $n=2000$. The regression parameters show comparatively stable performance throughout. Panel~(B) displays the empirical coverage probabilities of nominal 95\% confidence intervals. Consistent with the calibration results, the Poisson model achieves coverage probabilities close to the nominal level for all parameters, whereas the Binomial model exhibits undercoverage for the index parameters in moderate sample sizes due to variance underestimation. Nevertheless, the coverage probabilities steadily approach the target level as the sample size increases, demonstrating the asymptotic validity of the proposed inference procedure and highlighting the greater finite-sample difficulty associated with estimating nonlinear index parameters in logistic GPLSIMs.

%Recovery of the Nonlinear Link Function
% Figure~\ref{fig:estimated_link_function} and
Figure~\ref{fig:estimated_link_function_full} displays the estimated link function along with 95\% confidence bands for both Poisson and Binomial models across different sample sizes. Overall, the proposed estimator successfully captures the underlying nonlinear pattern of the true link function $\vp(u)=\sin(2u)$ in all settings. 
In the Poisson model (left panels), the estimated curves closely track the true function even at $n=400$, with only minor deviations near the boundaries. As the sample size increases to $n=800$ and $n=2000$, the estimation becomes nearly indistinguishable from the truth, and the confidence bands tighten substantially, indicating strong efficiency and consistency.
In contrast, the Binomial model (right panels) exhibits greater variability, particularly at $n=400$, where noticeable deviations from the true curve and wider confidence bands are observed, especially near the boundaries and regions with sparse data. However, as the sample size increases, the estimator improves markedly. At $n=800$, the overall shape is well recovered despite some remaining variability, and by $n=2000$, the estimated curve aligns closely with the true function, with significantly narrower confidence bands.
Across both models, the reduction in band width and improved alignment with the true function as $n$ increases provide clear evidence of consistency. The results also highlight that, while the Binomial setting is inherently more challenging, the proposed estimator remains robust and achieves accurate recovery of the link function with sufficient sample size.

% Figure~\ref{fig:estimated_link_function} shows the estimated link function with 95\% confidence bands. The proposed estimator accurately recovers $\vp(u)=\sin(2u)$ across both models. Deviations are minor and primarily occur near the boundaries, where data are sparse. Increasing the sample size from $n=800$ to $n=2000$ yields visibly narrower confidence bands and reduced variability, confirming consistency and improved precision. The pattern is consistent across both Poisson and Binomial settings.

% Figure~\ref{fig:estimated_link_function} presents the estimated nonlinear function $\vp(u)$ along with pointwise 95\% confidence bands. The results demonstrate that the proposed estimator is capable of accurately recovering the true nonlinear function across both Poisson and Binomial models. The estimated curves closely track the true function $\sin(2u)$ over the entire domain, with only minor deviations near the boundaries. These boundary effects are expected due to reduced data density in extreme regions of the index. A clear improvement is observed as the sample size increases from $n=800$ to $n=2000$, with narrower confidence bands and reduced variability. This reflects the increased precision of the spline estimator and confirms the expected convergence behavior. The overall pattern is consistent across both outcome types, indicating robustness of the method to different distributional assumptions. This evidence supports that uncertainty shrinks substantially with larger samples, while the mean estimate remains well-centered around the truth.

%Overall Summary
Overall, the proposed estimator achieves accurate parameter estimation, reliable recovery of nonlinear structure, and valid inference. It consistently outperforms the competing method in efficiency and stability, with gains most pronounced for index parameters and in smaller samples. The variance study further highlights the importance of correction methods for achieving nominal coverage in finite samples.

% The simulation studies demonstrate that the proposed penalized spline estimator provides accurate estimation, reliable inference, and robust recovery of nonlinear effects in generalized partially linear single-index model settings. The method performs consistently across different outcome types and sample sizes, with clear advantages in efficiency and stability over existing approaches. Furthermore, the variance study underscores the importance of robust variance estimation in finite samples, while confirming that the proposed framework achieves valid inference asymptotically.


% \subsection*{Figure 2: Power Curves of Spline Wald Tests}

% \begin{figure}[ht]
% \centering

% %----- First row: alpha parameters -----%
% \begin{subfigure}[t]{0.32\textwidth}
%     \centering
%     \includegraphics[width=\textwidth]{par1Poisson_PowerAnalysis.pdf}
%     \caption{$\alpha_{10}$}
%     \label{fig:power_alpha1}
% \end{subfigure}
% \hfill
% \begin{subfigure}[t]{0.32\textwidth}
%     \centering
%     \includegraphics[width=\textwidth]{par2Poisson_PowerAnalysis.pdf}
%     \caption{$\alpha_{20}$}
%     \label{fig:power_alpha2}
% \end{subfigure}
% \hfill
% \begin{subfigure}[t]{0.32\textwidth}
%     \centering
%     \includegraphics[width=\textwidth]{par3Poisson_PowerAnalysis.pdf}
%     \caption{$\alpha_{30}$}
%     \label{fig:power_alpha3}
% \end{subfigure}

% \vspace{0.4cm}

% %----- Second row: beta parameters -----%
% \begin{subfigure}[t]{0.32\textwidth}
%     \centering
%     \includegraphics[width=\textwidth]{par4Poisson_PowerAnalysis.pdf}
%     \caption{$\beta_{10}$}
%     \label{fig:power_beta1}
% \end{subfigure}
% \hfill
% \begin{subfigure}[t]{0.32\textwidth}
%     \centering
%     \includegraphics[width=\textwidth]{par5Poisson_PowerAnalysis.pdf}
%     \caption{$\beta_{20}$}
%     \label{fig:power_beta2}
% \end{subfigure}

% \caption{\textbf{Power curves of spline-based Wald tests for the Poisson model.}
% Each panel displays the empirical power for testing the null hypotheses
% $H_0: \alpha_{10}$, $H_0: \alpha_{20}$, $H_0: \alpha_{30}$, $H_0: \beta_{10}$, and
% $H_0: \beta_{20}$, respectively. The vertical red line indicates the true parameter value, while the horizontal dashed line represents the nominal 5\% significance level. Green and blue curves correspond to sample sizes $n = 800$ and $n = 2000$, respectively. Results are based on 2000 Monte Carlo simulation replications using the conditional expected information variance estimator.}
% \label{fig:power}
% \end{figure}

\newpage

\subsection{Real Data Analysis}

We illustrate the proposed generalized partially linear single-index model (GPLSIM) using data from the National Health and Nutrition Examination Survey (NHANES), a nationally representative cross-sectional study designed to assess the health and nutritional status of adults and children in the United States \cite{nhanes}. The data are accessed via the \texttt{NHANES} R package \cite{nhanesR}, which provides a curated subset suitable for statistical analysis. We consider a binary diabetes outcome, \textit{Diabetes}. The single-index component includes body mass index (\textit{BMI}), systolic blood pressure (\textit{BPSysAve}), and directly measured HDL cholesterol (\textit{DirectChol}), while the linear component includes age, gender (\textit{Gender}), and race. To improve stability and interpretability, the original \textit{Hispanic} and \textit{Mexican} race categories are combined into a single \textit{Hispanic} category because of their relatively small sample sizes, and the \textit{White} group is used as the reference category in the regression analysis. Continuous predictors are standardized prior to model fitting. We compare three approaches: a parametric generalized linear model (GLM) \cite{mccullagh1989glm}, a nonparametric generalized additive model (GAM) \cite{hastie1990gam}, and the proposed GPLSIM (Lu GPLSIM).

% \begin{table}[ht!]
% \centering
% \caption{Summary of the NHANES working dataset used in the GPLSIM analysis.}
% \label{tab:summary_realdata}
% \small
% \begin{tabular}{llccc}
% \toprule
% \textbf{Role} & \textbf{Variable} &
% \textbf{Mean (SD)} &
% \textbf{Median} &
% \textbf{Range / Count (\%)} \\
% \midrule

% Outcome
% & Diabetes (Yes)
% & --
% & --
% & 674 (8.50\%) \\

% \midrule

% \multirow{3}{*}{Single-index variables}
% & BMI
% & 27.73 (6.94)
% & 26.73
% & [12.90, 81.25] \\

% & BPSysAve
% & 118.30 (17.06)
% & 116.0
% & [76.0, 226.0] \\

% & DirectChol
% & 1.37 (0.40)
% & 1.29
% & [0.39, 4.03] \\

% \midrule

% \multirow{5}{*}{Linear covariates}
% & Age
% & 41.45 (19.86)
% & 41.0
% & [8, 80] \\

% & Gender (Male)
% & --
% & --
% & 3958 (49.90\%) \\

% & Race: White
% & --
% & --
% & 5229 (65.93\%) \\

% & Race: Black
% & --
% & --
% & 878 (11.07\%) \\

% & Race: Hispanic
% & --
% & --
% & 1225 (15.45\%) \\

% & Race: Other
% & --
% & --
% & 599 (7.55\%) \\

% \bottomrule
% \end{tabular}

% \vspace{0.15cm}
% \begin{minipage}{0.95\textwidth}
% \footnotesize
% \textit{Note:} The analysis includes $n=7{,}931$ participants. The binary response variable is \textit{Diabetes}. The single-index component consists of body mass index (\textit{BMI}), average systolic blood pressure (\textit{BPSysAve}), and directly measured HDL cholesterol (\textit{DirectChol}). The linear component includes age, gender, and race/ethnicity, where the Hispanic category combines the original Mexican American and Other Hispanic groups. Continuous variables are summarized by mean (standard deviation), median, and range, whereas categorical variables are summarized by counts and percentages. Continuous predictors are standardized before model fitting.
% \end{minipage}
% \end{table}

\begin{table}[ht!]
\centering
\caption{Comparison of the estimated effects of the single-index variables under the GLM, GAM, and proposed GPLSIM (Lu GPLSIM).}
\label{tab:singleindex_compare}
\small
\begin{tabular}{p{2.2cm}p{3.8cm}p{3.8cm}p{3.8cm}}
\toprule
\textbf{Model / Quantity}
&
\textbf{BMI}
&
\textbf{SBP}
&
\textbf{HDL cholesterol (DirectChol)}
\\
\midrule

\multicolumn{4}{l}{\textbf{GLM}}\\

Estimate ($p$-value) &
OR = 1.76 ($^{***}$) &
OR = 1.07 (NS) &
OR = 0.73 ($^{***}$) \\

Interpretation &
Significant positive linear effect &
No significant linear association &
Significant protective linear effect \\

\addlinespace

\multicolumn{4}{l}{\textbf{GAM}}\\

Estimate ($p$-value) &
EDF = 2.382 ($^{***}$) &
EDF = 3.591 (NS) &
EDF = 1.002 ($^{***}$) \\

Interpretation &
Significant nonlinear smooth effect &
No significant smooth effect &
Significant approximately linear smooth effect \\

\addlinespace

\multicolumn{4}{l}{\textbf{Lu GPLSIM}}\\

Estimate ($p$-value) &
$\hat{\alpha}_1 = 0.839$ ($^{***}$) &
$\hat{\alpha}_2 = 0.135$ ($^{**}$) &
$\hat{\alpha}_3 = -0.527$ ($^{***}$) \\

Interpretation &
Largest positive contributor to the latent index &
Small positive contributor to the latent index &
Large negative contributor to the latent index \\

\bottomrule
\end{tabular}

\vspace{0.15cm}

\begin{minipage}{0.96\textwidth}
\footnotesize
\textit{Note:}
For the GLM, estimates are reported as odds ratios (ORs). For the GAM, estimates are summarized by the effective degrees of freedom (EDF) of the smooth function, where EDF close to one indicates an approximately linear relationship and larger EDF indicates increasing nonlinearity. For the proposed GPLSIM, the estimated coefficients quantify the relative contribution of each variable to the latent single index rather than individual log-odds effects. Significance codes are defined as $^{***}p<0.001$, $^{**}p<0.01$, $^{*}p<0.05$, and NS: not statistically significant ($p\ge0.05$). Across all three approaches, BMI consistently emerges as the strongest positive contributor to diabetes risk, HDL cholesterol exhibits a protective association, and systolic blood pressure contributes comparatively little after adjustment for the remaining covariates.
\end{minipage}

\end{table}


% Compare the relative importance and direction of each single-index variable (i.e., BMI, SBP, and DirectChol) with the corresponding estimation results from either GLM or GAM.
The estimated effects of the three single-index variables are remarkably consistent across the GLM, GAM, and the proposed GPLSIM, although the three models represent these effects differently. The GLM summarizes each predictor through a constant odds ratio, the GAM models each predictor by a separate smooth function characterized by its effective degrees of freedom (EDF), and the proposed GPLSIM combines the variables into a common latent metabolic index whose coefficients quantify their relative contributions and directions rather than individual odds ratios.

For body mass index (BMI), the GLM indicates a strong positive association with diabetes, with an estimated odds ratio of 1.76 ($p<0.001$). The GAM likewise identifies BMI as a highly significant nonlinear predictor ($\mathrm{EDF}=2.382$, $p<0.001$), indicating that its effect cannot be adequately represented by a simple linear relationship. Consistent with these findings, the proposed GPLSIM assigns BMI the largest positive index coefficient ($\hat{\alpha}_1=0.839$, $p<0.001$), indicating that BMI is the dominant contributor to the latent metabolic index. Thus, all three approaches consistently identify BMI as the strongest metabolic risk factor for diabetes.

For systolic blood pressure (SBP), the GLM estimates an odds ratio of 1.07, which is not statistically significant ($p=0.136$), suggesting little independent association with diabetes after adjusting for the remaining covariates. Similarly, the GAM reports a nonsignificant smooth effect ($\mathrm{EDF}=3.591$, $p=0.147$), providing no evidence of a meaningful nonlinear relationship. In contrast, the GPLSIM assigns SBP a small positive index coefficient ($\hat{\alpha}_2=0.135$, $p<0.01$). Although this coefficient is statistically significant, its magnitude is substantially smaller than that of BMI and HDL cholesterol, indicating that SBP contributes relatively little to the construction of the latent metabolic index. Therefore, despite its statistical significance within the index formulation, SBP appears to play only a minor role in the overall metabolic risk profile.

For directly measured HDL cholesterol (DirectChol), the GLM estimates an odds ratio of 0.73 ($p<0.001$), indicating a significant protective association with diabetes. The GAM likewise identifies a highly significant smooth effect with an effective degree of freedom close to one ($\mathrm{EDF}=1.002$, $p<0.001$), implying an approximately linear protective relationship. The proposed GPLSIM similarly assigns HDL cholesterol a large negative index coefficient ($\hat{\alpha}_3=-0.527$, $p<0.001$), indicating that increasing HDL levels reduce the latent metabolic index. Assuming a monotone increasing link function, a lower latent index corresponds to a lower probability of diabetes, consistent with the protective effect observed in both the GLM and GAM. Moreover, the magnitude of the GPLSIM coefficient suggests that HDL cholesterol is the second most influential component of the latent metabolic index after BMI.

Overall, the three modeling approaches produce highly consistent substantive conclusions regarding the direction and relative importance of the metabolic risk factors. BMI exhibits the strongest positive contribution to diabetes risk, HDL cholesterol exhibits a substantial protective contribution, and systolic blood pressure contributes comparatively little after accounting for the remaining variables. The principal distinction among the models lies in their representation of these effects: the GLM expresses them through constant odds ratios, the GAM through flexible smooth functions, and the proposed GPLSIM through a latent single index whose coefficients quantify the relative importance of the metabolic variables.

The ability to summarize multiple correlated metabolic risk factors through a single latent index is particularly advantageous in practice. Rather than interpreting several individual nonlinear effects separately, clinicians and public health researchers can evaluate an individual's overall metabolic status using a single composite score that integrates the joint contributions of BMI, systolic blood pressure, and HDL cholesterol. This lower-dimensional representation facilitates risk stratification, visualization, and prediction while mitigating multicollinearity among correlated biomarkers. Furthermore, because the latent index is estimated directly from the data, it provides an interpretable summary measure of overall metabolic risk that can be incorporated into downstream analyses or clinical decision-making without requiring separate assessment of each nonlinear covariate effect. Consequently, the proposed GPLSIM offers a parsimonious yet flexible framework that preserves the essential nonlinear structure of the metabolic variables while yielding a clinically meaningful summary of diabetes risk.


\begin{table}[ht!]
\centering
\caption{Comparison of the estimated coefficients for the linear covariates across the GLM, GAM, and GPLSIM.}
\label{tab:linear_covariates}
\small
\begin{tabular}{lccc}
\toprule
\textbf{Covariate}
& \textbf{GLM Estimate}
& \textbf{GAM Estimate}
& \textbf{GPLSIM Estimate} \\
\midrule

Age (standardized)
&
1.216***
&
1.206***
&
1.212*** \\

Male (vs. Female)
&
0.177
&
0.185
&
0.164 \\

Black (vs. White)
&
0.778***
&
0.785***
&
0.822*** \\

Hispanic (vs. White)
&
0.664***
&
0.655***
&
0.654*** \\

Other (vs. White)
&
0.842***
&
0.847***
&
0.854*** \\

\bottomrule
\end{tabular}

\vspace{0.15cm}
\begin{minipage}{0.96\textwidth}
\footnotesize
\textit{Note:}
The proposed GPLSIM models age, gender, and race/ethnicity through a linear component, so their coefficients retain the same log-odds interpretation as those from a conventional logistic regression model and may be exponentiated to obtain adjusted odds ratios. White serves as the reference category for race/ethnicity, and the Hispanic category combines the original Mexican American and Other Hispanic groups. The estimated GPLSIM coefficients ($\hat{\alpha}$) are directly comparable with the corresponding GLM and GAM coefficients, facilitating assessment of the consistency of the linear effects across the three competing models. Significance codes: $^{***}p<0.001$.
\end{minipage}
\end{table}



% Compare the estimated coefficients of the linear covariates with those results from either GLM or GAM.
Unlike the single-index variables, whose effects are summarized through a latent nonlinear index, the linear covariates in the proposed GPLSIM retain the same interpretation as those in a conventional logistic regression model. Consequently, their estimated coefficients can be directly compared with those obtained from the GLM and the parametric component of the GAM. As shown in Table~\ref{tab:linear_covariates}, the estimated coefficients are remarkably consistent across the three modeling approaches. Age exhibits a strong positive association with diabetes risk in all models, with estimated coefficients ranging from 1.206 to 1.216, indicating that older individuals are substantially more likely to have diabetes after adjusting for the remaining covariates. The estimated gender effect is also highly stable, with positive coefficients between 0.164 and 0.185, suggesting that males tend to have a higher probability of diabetes than females, although the corresponding effect is not statistically significant at the conventional 5\% level.

The estimated race effects likewise demonstrate excellent agreement across the three models. Using White as the reference category, Black individuals exhibit a significantly higher risk of diabetes, with estimated coefficients ranging from 0.778 to 0.822. Similarly, the combined Hispanic category, which merges the original Mexican American and Other Hispanic groups, shows a significantly elevated diabetes risk relative to White individuals, with estimated coefficients around 0.65 across all models. The Other race category also exhibits a consistently positive association with diabetes, with estimated coefficients between 0.842 and 0.854. The close numerical agreement of these estimates demonstrates that the proposed GPLSIM preserves the interpretation of the demographic covariates while allowing greater modeling flexibility for the metabolic variables represented through the single-index component.

Because age, gender, and race enter the GPLSIM through a linear component, their coefficients retain the standard log-odds interpretation and can be exponentiated to obtain adjusted odds ratios exactly as in a conventional logistic regression model. This property distinguishes the GPLSIM from a fully nonparametric additive model, where nonlinear smooth functions are estimated for all continuous variables and direct odds-ratio interpretations are generally unavailable. Consequently, the proposed GPLSIM combines the interpretability of classical logistic regression for demographic covariates with the flexibility of a nonlinear latent-index representation for the correlated metabolic predictors, yielding a parsimonious yet scientifically interpretable framework for diabetes risk assessment. The high level of agreement between the GPLSIM and the corresponding GLM and GAM estimates further indicates that incorporating a nonlinear latent metabolic index does not materially alter inference for the linear demographic covariates while substantially enriching the representation of the correlated metabolic risk factors.





\begin{figure}[ht!]
\centering

\includegraphics[width=0.9\textwidth]{Binomial_RealData_260619.png}

\caption{
Estimated link functions along the single-index direction for the binary outcome of diabetes (\textit{Diabetes}). The left panel compares the GLM (black) with the Lu GPLSIM estimator (blue), while the right panel compares the GLM (black) with the Yu GPLSIM estimator (red). Dashed curves denote pointwise 95\% confidence intervals for the estimated nonparametric link functions. % Both GPLSIM approaches reveal departures from the linear logit relationship assumed by the GLM, particularly in the lower and middle regions of the single-index domain, while producing similar monotone increasing trends overall.
}
\label{fig:realdata_link}
\end{figure}

% Figure~\ref{fig:realdata_link} presents the estimated link functions along the single-index direction. For the Poisson outcome (top panels), the GLM imposes a linear structure that fails to capture the evident nonlinear relationship between predictors and the outcome. Both the GAM and GPLSIM reveal a clear nonlinear, S-shaped pattern, particularly in the central region of the index. The GPLSIM closely follows the nonlinear trend identified by the GAM while producing a smoother and more stable estimate, especially near the boundaries where variability tends to increase. For the binary outcome (bottom panels), the deviation from linearity is more moderate. The GAM suggests mild curvature, and the GPLSIM captures a similar nonlinear pattern, with slightly stronger curvature in the upper range of the index. Although uncertainty increases near the boundaries due to limited data support, the overall agreement between the GAM and GPLSIM indicates the presence of modest but non-negligible nonlinear effects.

% These results demonstrate that the GPLSIM provides a parsimonious and interpretable framework for modeling nonlinear relationships. Compared to the GLM, it offers substantially improved flexibility, while avoiding the excessive variability that can arise in fully nonparametric approaches such as the GAM.

Figure~\ref{fig:realdata_link} presents the estimated link functions along the single-index direction for the binary diabetes outcome. The left panel compares the GLM and the Lu GPLSIM estimator, while the right panel compares the GLM and the Yu GPLSIM estimator. The GLM imposes a strictly linear relationship between the single index and the log-odds of diabetes, whereas both GPLSIM approaches estimate the link function nonparametrically. Overall, the two GPLSIM estimators reveal a nonlinear monotone increasing relationship between the latent metabolic index and diabetes risk that is not fully captured by the linear GLM.

The Lu GPLSIM estimate exhibits a smooth S-shaped pattern, with a relatively gradual increase for smaller index values, followed by a steeper rise in the middle region before approaching an approximately linear trend for larger index values. The estimated confidence bands remain relatively narrow across most of the index domain, indicating stable estimation except near the boundaries where observations are sparse. In comparison, the Yu GPLSIM produces a qualitatively similar nonlinear pattern but displays stronger curvature in the lower-index region and somewhat wider confidence intervals, reflecting greater variability in estimating the link function near the extremes of the observed index values. Despite these local differences, both GPLSIM approaches consistently suggest that the association between the latent metabolic score and diabetes risk deviates from the linear relationship assumed by the GLM.

Overall, these results demonstrate the advantage of generalized partially linear single-index models for modeling complex metabolic effects on diabetes risk. By combining multiple correlated biomarkers into a single latent metabolic index while estimating its nonlinear effect nonparametrically, both GPLSIM methods provide substantially greater flexibility than the conventional GLM. Moreover, the close agreement between the Lu and Yu GPLSIM estimators indicates that the inferred nonlinear relationship is robust across estimation procedures, while the Lu GPLSIM yields a smoother estimate with narrower confidence bands over much of the index range.



% @misc{nhanes,
%   author = {{Centers for Disease Control and Prevention (CDC)}},
%   title = {National Health and Nutrition Examination Survey (NHANES)},
%   year = {2020},
%   note = {https://www.cdc.gov/nchs/nhanes/}
% }

% @manual{nhanesR,
%   title        = {NHANES: Data from the US National Health and Nutrition Examination Study},
%   author       = {Goldberg, Rachel and others},
%   year         = {2020},
%   note         = {R package version 2.1.0},
%   url          = {https://CRAN.R-project.org/package=NHANES}
% }

% @book{mccullagh1989glm,
%   author = {McCullagh, P. and Nelder, J. A.},
%   title = {Generalized Linear Models},
%   edition = {2nd},
%   year = {1989},
%   publisher = {Chapman \& Hall}
% }

% @book{hastie1990gam,
%   author = {Hastie, T. and Tibshirani, R.},
%   title = {Generalized Additive Models},
%   year = {1990},
%   publisher = {Chapman \& Hall}
% }

\newpage
